\section{Two test functions}\label{cf:sec:test}

\begin{proposition}[the two ends]\label{cf:prop:ends}
\begin{enumerate}
\item[(a)] (cited) For $\theta\ge1$, unconditionally, $u_\theta(v)=\sum_\rho P_{\theta-\sigma_\rho}(v-\gamma_\rho)>0$; in particular $\mathfrak U_1=\sum_\rho P_{1-\sigma_\rho}(\cdot-\gamma_\rho)$ is a positive mixture of Cauchy laws centred at the ordinates, with scales in $(0,1)$. This is Theorem~\ref{lo:thm:density} and Corollary~\ref{lv:cor:uncond}, and it holds for $\zeta$ because $\zeta(1+it)\ne0$.
\item[(b)] (proved) As $\theta\downarrow\frac12$, in the sense of distributions, $u_\theta\to\mathfrak U_{1/2}=\sum_{\sigma_\rho=1/2}\delta_{\gamma_\rho}$. Zeros off the line are invisible in $\mathfrak U_{1/2}$: a pair $\sigma_\rho$, $1-\sigma_\rho$ contributes $Q_{\theta-\sigma_\rho}+Q_{\theta-1+\sigma_\rho}\to Q_{1/2-\sigma_\rho}+Q_{\sigma_\rho-1/2}=0$.
\end{enumerate}
\end{proposition}

\begin{proof}
(b) For a zero on the line $P_{\theta-1/2}\to\delta_0$. For a pair, $\theta-\sigma_\rho$ and $\theta-1+\sigma_\rho$ tend to $\frac12-\sigma_\rho\ne0$ and $\sigma_\rho-\frac12\ne0$, and $a\mapsto Q_a$ is continuous from $\mathbb R\setminus\{0\}$ to $L^1(\mathbb R)$, so the pair terms converge to $-P_{\sigma_\rho-1/2}+P_{\sigma_\rho-1/2}=0$. The far terms are bounded by $\lvert\gamma_\rho\rvert^{-2}$ uniformly for $\theta\in[\frac12,1]$, which gives dominated convergence on compact sets.
\end{proof}

In the variable $x=v^2$ this is the statement that $U_{1/2}=\sum_\gamma\delta_{\gamma^2}$ is the discrete measure of the Riemann hypothesis programme.

\begin{resultpar}\label{cf:res:dhzeros}
All computations of this chapter use $\zeta$ and the Davenport and Heilbronn function, which has no Euler product and has zeros off the critical line. With $\kappa=(\sqrt{10-2\sqrt5}-2)/(\sqrt5-1)=0.28407904384\ldots$ it is \cite{DH36,Titchmarsh}
\[
f(s)=5^{-s}\Bigl[\zeta(s,\tfrac15)+\kappa\,\zeta(s,\tfrac25)-\kappa\,\zeta(s,\tfrac35)-\zeta(s,\tfrac45)\Bigr],
\]
and the completed function $\Lambda(s)=(5/\pi)^{s/2}\Gamma(\frac{s+1}2)f(s)$ is entire and satisfies $\Lambda(s)=\Lambda(1-s)$; the identity is that of Davenport and Heilbronn \cite{DH36} (see \cite{Titchmarsh}), and as a diagnostic the relative difference $\lvert\Lambda(s)-\Lambda(1-s)\rvert/\lvert\Lambda(s)\rvert$ is at most $1.1\times10^{-87}$ at four points in ball arithmetic. The function is treated in Section~\ref{ht:sec:dhtest}; here it is the control. The zeros of $f$ off the line below height $200$ were found by Spira \cite{Spira94}; refined by Newton's method in ball arithmetic (\texttt{code/cf\_dhzeros.py}, \texttt{data/cf\_dhzeros.json}) they are
\begin{gather*}
0.8085171825+85.6993484854\,i,\qquad 0.6508300806+114.1633427308\,i,\\
0.5743560505+166.4793059132\,i,\qquad 0.7242576946+176.7024612429\,i,
\end{gather*}
together with their reflections in the critical line, each certified to be the only zero in a square of half width $10^{-16}$ about the centre, so that all ten printed decimals are correct. These are all the zeros off the line below height $200$ (computed, certified in ball arithmetic): $\Lambda$ has $130$ zeros with $0<\gamma<200$, counted by the argument principle on a rectangle, and $122$ sign changes of $\Lambda(\frac12+it)$ on the segment, with certified signs on a grid of step $\frac1{50}$; the difference $8=2\cdot4$ is the number of off line zeros and their reflections, and equality means that no zero on the line is missed or multiple and that no zero off the line other than these four and their reflections exists.
\end{resultpar}

For $\Lambda$ the density $u_\theta=\pi^{-1}\Re\Lambda'/\Lambda(\theta+iv)$ plays the role of $u_\theta$ for $\xi$. The Hadamard sum~\eqref{cf:eq:mg}, the transport identity of Proposition~\ref{cf:prop:transport} and the deconvolution criterion with the base line $\Re s=1$ replaced by a line to the right of all zeros apply to it ($\Re s=2$ works, since $\lvert f(s)-1\rvert\le\zeta(2)-1<1$ there, as $\lvert\kappa\rvert<1$); Proposition~\ref{cf:prop:ends}(a) does not apply, since $f$ has infinitely many zeros in $\Re s>1$ \cite{DH36}. Proposition~\ref{cf:prop:barrier} makes this precise. The principle of such a test is that of Section~\ref{ht:sec:dhtest}: a method that detects zeros off the line of $\xi$ must detect them for $\Lambda$.

For $\zeta$ the densities $u_\theta$ on the window $[60,120]$ stay positive at every level computed (the zeros of $\zeta$ at these heights are on the line \cite{PT21}) and sharpen into spikes at the ordinates as $\theta$ decreases, which is the transport identity with $R_\theta=0$. For $\Lambda$ the densities at $\theta=1$ and $\theta=0.9$ are positive on this window (computed; the zeros of $f$ in the window have $\sigma_\rho\le0.8085$). They turn negative near $v=85.70$ once $\theta<0.8085$ and near $v=114.16$ once $\theta<0.6508$: at the ordinates the minima of $u_\theta$ over $[60,120]$ are $-36.8$ at $\theta=0.8$ and $-4.96$ at $\theta=0.6$ (computed, \texttt{code/cf\_xi.py}; the grid values $-36.6$ and $-4.6$ of the note are grid minima). This is the signed defect of Proposition~\ref{cf:prop:transport}, seen directly.

\section{A particle filter along the Cauchy flow}\label{cf:sec:pf}

\subsection{The state space model}\label{cf:sec:model}

Fix a window of heights. The parameter is the local configuration $z=\{(\gamma_j,\sigma_j)\}_{j=1}^K$ of the zeros with ordinates in a slightly larger window, with $K$ fixed by the argument principle (a contour integral of $\xi'/\xi$, computed once). The observation at level $\theta_k$ is the vector $y_k=(u_{\theta_k}(v_i))_i$ on a grid, computed to high precision. Equation~\eqref{cf:eq:mg} gives the observation equation
\[
y_k(v)=\sum_{j=1}^KQ_{\theta_k-\sigma_j}(v-\gamma_j)+b_k(v)+\varepsilon_k(v),
\]
where $b_k$ is the contribution of the zeros outside the window. For a distant zero the pair terms equal $2(\theta-\frac12)/(\pi(v-\gamma)^2)$ to leading order, so we model $b_k(v)=(\theta_k-\frac12)(c_0+c_1w+c_2w^2+c_3w^3)+(\theta_k-\frac12)^3c_4$ with $w$ the centred height. The coefficients $c$ enter linearly and are integrated out in closed form under a Gaussian prior, by the Woodbury identity, so the particles carry only $(\gamma_j,\sigma_j)$. The error $\varepsilon$ is a tolerance, not noise: the data are deterministic, and $\varepsilon=0.01$ absorbs the background approximation. Every statement below about posteriors, information and coverage refers to this Gaussian tolerance likelihood.

The levels $1=\theta_1>\theta_2>\dots>\theta_n$ play the role of time. With the background coefficients shared by the levels and integrated out jointly, the filtering distributions are
\[
\pi_k(z)\;\propto\;\pi_0(z)\,p(y_1,\dots,y_k\mid z),
\]
the Feynman Kac flow of \cite{PSH626} with potentials $G_k=p(y_k\mid y_1,\dots,y_{k-1},z)$ and Markov moves invariant for $\pi_k$. The prior is uniform: ordinates uniform in the window and every $\sigma_j$ uniform on $(0,1)$, so that the functional equation is not imposed on the particles (Section~\ref{cf:sec:detect}).

\subsection{Algorithm}

\begin{enumerate}
\item Draw $N$ particles from the prior.
\item For $k=1,\dots,n$: move from $\pi_{k-1}$ to $\pi_k$ through the tempered targets $\pi_{k-1}^{1-\lambda}\pi_k^{\lambda}$, choosing each increment of $\lambda$ by bisection so that the effective sample size of the incremental weights is $N/2$; resample systematically; rejuvenate by componentwise random walk Metropolis moves with scales set from the current particle cloud.
\item Report the posterior of every $\sigma_j$, and of $\delta^\ast=\max_j\lvert\sigma_j-\frac12\rvert$ over the ordinates in the interior of the window.
\end{enumerate}
This is the sequential Monte Carlo sampler of Chopin \cite{Chopin02} and Del Moral, Doucet and Jasra \cite{DDJ06}, with the Cauchy level as time; the book's account of sequential Monte Carlo for the explicit formula is Section~\ref{es:sec:adaptive}, and the single level Bayesian inversion is Section~\ref{lo:sec:bayes}, with its path of Section~\ref{lo:sec:path}.

\subsection{Results for $\zeta$}

The runs of this section and the next are the note's runs (computed by the author, not recomputed for the book). We used $N=1000$ particles, the window $[78,94]$ with $401$ grid points, $K=10$ zeros in $[74,98]$, and levels $\theta\in\{1,0.9,0.8,0.7,0.65,0.6\}$. Table~\ref{cf:tab:zeta} gives the posterior for the zeros in the data window and Table~\ref{cf:tab:flow} the posterior of $\delta^\ast$ as levels are added. The ordinates are recovered to within $10^{-4}$ and every $\sigma_j$ is pinned at $\frac12$ to within $4\times10^{-4}$. The bound on $\delta^\ast$ falls by a factor of fourteen between $\theta=1$ and $\theta=0.6$. Each $\sigma_j$ is a free parameter, so the information is first order, $\int(\partial_aP_a)^2\,dx=1/(4\pi a^3)$, and $(\sum_la_l^{-3}/a_1^{-3})^{1/2}=13.6$ for these levels (Section~\ref{cf:sec:detect}); Proposition~\ref{cf:prop:detect}, which concerns a symmetric pair, would give $7.8$. Zeros just outside the data window are poorly identified, as they should be; they act as nuisance parameters.

\begin{table}[h]
\centering
\begin{tabular}{ccc}
\toprule
posterior median $\gamma$ & posterior median $\sigma_j$ & $95\%$ interval for $\sigma_j$\\
\midrule
$79.3373$ & $0.49999$ & $[0.49964,\,0.50036]$\\
$82.9104$ & $0.50000$ & $[0.49966,\,0.50035]$\\
$84.7355$ & $0.50002$ & $[0.49967,\,0.50038]$\\
$87.4253$ & $0.49999$ & $[0.49964,\,0.50034]$\\
$88.8091$ & $0.49999$ & $[0.49961,\,0.50036]$\\
$92.4919$ & $0.50001$ & $[0.49967,\,0.50037]$\\
\bottomrule
\end{tabular}
\caption{Particle filter posterior for the zeros of $\zeta$ in the data window $[78,94]$ after all levels down to $\theta=0.6$ (computed by the author). The true ordinates are $79.3374$, $82.9104$, $84.7355$, $87.4253$, $88.8091$ and $92.4919$.}
\label{cf:tab:zeta}
\end{table}

\begin{table}[h]
\centering
\begin{tabular}{lcccccc}
\toprule
levels absorbed down to $\theta=$ & $1$ & $0.9$ & $0.8$ & $0.7$ & $0.65$ & $0.6$\\
\midrule
median of $\delta^\ast$ ($\times10^{-3}$) & $3.65$ & $2.15$ & $1.31$ & $0.73$ & $0.46$ & $0.26$\\
$95\%$ quantile of $\delta^\ast$ ($\times10^{-3}$) & $6.43$ & $3.69$ & $2.33$ & $1.34$ & $0.85$ & $0.46$\\
likelihood evaluations ($\times10^{6}$) & $1.86$ & $2.46$ & $2.95$ & $3.43$ & $3.79$ & $4.15$\\
\bottomrule
\end{tabular}
\caption{The backward Cauchy flow as a filter: the posterior of $\delta^\ast=\max\lvert\sigma_j-\frac12\rvert$ over the zeros of $\zeta$ in $[80,92]$ as levels are added, with cumulative cost (computed by the author). The bounds measure the displacement of one unpaired zero from the line, which the functional equation forbids.}
\label{cf:tab:flow}
\end{table}

For a pair $\frac12\pm\delta$ the same data bound $\delta$ only at about $4.4\times10^{-3}$ (one sided $95\%$, second order Gaussian approximation, six levels; $0.0345$ after level one alone), ten times weaker than the bound on an unpaired displacement (computed, \texttt{code/cf\_info.py}). A rerun with the pair parametrisation of Section~\ref{es:sec:singular} (Lemma~\ref{es:lem:even}, Theorem~\ref{es:thm:singular}), each ordinate on the line or a pair $\frac12\pm\delta$ with $\eta=\delta^2\ge0$, has not been done.

The cost is the obstacle. The run used $4.15\times10^6$ likelihood evaluations for a window of length $16$ containing six zeros. For the Davenport and Heilbronn window, with seventeen zeros and a uniform prior, the tempering stalled at $\lambda\approx0.003$ after the first several hundred thousand evaluations: the posterior is so concentrated relative to the prior that the bridge needs thousands of steps. Optimising this is the job of GBC (Section~\ref{cf:sec:gbc}).

\subsection{How much information each level carries}\label{cf:sec:detect}

\begin{proposition}[detectability]\label{cf:prop:detect}
Let $a=\theta-\frac12>\delta>0$. At the centre $v=\gamma$, a pair of zeros at $\frac12\pm\delta+i\gamma$ exceeds a double zero on the line at $\frac12+i\gamma$ by
\[
\Delta(a,\delta)=\frac1\pi\Bigl(\frac1{a-\delta}+\frac1{a+\delta}\Bigr)-\frac2{\pi a}=\frac{2\delta^2}{\pi a(a^2-\delta^2)}.
\]
In the parameter $\eta=\delta^2$ the pair differs from the double zero by $\eta\,\partial_a^2P_a+O(\eta^2)$, and
\[
\int_{\mathbb R}\bigl(\partial_a^2P_a(x)\bigr)^2dx=\frac3{4\pi a^5}.
\]
\end{proposition}

\begin{proof}
Direct computation, with $\partial_a^2P_a(x)=2a(a^2-3x^2)/\bigl(\pi(a^2+x^2)^3\bigr)$ and the scaling $P_a(x)=a^{-1}P_1(x/a)$. The integral was also checked by quadrature at $40$ digits (\texttt{code/cf\_info.py}).
\end{proof}

Two consequences concern a symmetric pair, the situation of Section~\ref{es:sec:singular}. First, the information about $\delta$ is second order: the Fisher information at $\delta=0$ vanishes, and one works with $\eta=\delta^2$. Second, the information grows like $a^{-5}$ as the level approaches the line. With observation error $\varepsilon$ on a grid of spacing $\Delta v$, the level $a$ resolves $\delta$ once roughly $\delta^4\gtrsim(4\pi/3)\,\varepsilon^2\Delta v\,a^5$ (heuristic, for the pair model). Coarse levels are well posed and uninformative; fine levels are informative and ill posed, which is the structure that sequential Monte Carlo handles well: the level is the filter's time, and each new level sharpens the likelihood of a pair in $\eta$ by a factor of order $((\theta_{k-1}-\frac12)/(\theta_k-\frac12))^5$. The filter of this section does not impose the pairing, so its information is first order, as follows.

\subsection{Information rates}

For a single zero with free real part $\sigma$ (no pairing imposed), $\partial_\sigma Q_{\theta-\sigma}=-\partial_aP_a$ and $\int(\partial_aP_a)^2dx=1/(4\pi a^3)$. With error $\varepsilon$ on a grid of spacing $\Delta v$, the Fisher information after levels $a_1,\dots,a_k$, with $a_l=\theta_l-\frac12$, is about $(4\pi\varepsilon^2\Delta v)^{-1}\sum_{l\le k}a_l^{-3}$, so the Bernstein and von Mises posterior standard deviation of the Gaussian tolerance model is
\[
\mathrm{sd}(\sigma_j)\approx\varepsilon\sqrt{4\pi\Delta v}\,\Bigl(\sum_{l\le k}a_l^{-3}\Bigr)^{-1/2}\qquad\text{(heuristic).}
\]
Table~\ref{cf:tab:info} compares this with the particle filter of this section. The agreement is within $8\%$ at the first level and exact at the last, so the filter matches the first order Fisher information of the model with free real parts; a deterministic tolerance likelihood gives no frequentist coverage, and the efficient information for the full multizero model with unknown background is not that of an isolated zero, so no more is claimed. The ratios of the median bound of Table~\ref{cf:tab:flow} to its value at level one (observed $1.70$ to $1.74$, $2.76$ to $2.79$, $4.80$ to $5.00$, $7.56$ to $7.93$, $13.98$ to $14.04$) follow the first order law, $(\sum_la_l^{-3}/a_1^{-3})^{1/2}=1.72$, $2.75$, $4.82$, $7.76$, $13.61$ for the levels $0.9$ to $0.6$, and not the law of a pair, $(\sum_la_l^{-5}/a_1^{-5})^{1/4}=1.42$, $2.03$, $3.27$, $4.79$, $7.77$ (computed, \texttt{code/cf\_info.py}).

\begin{table}[h]
\centering
\begin{tabular}{lcccccc}
\toprule
levels absorbed down to $\theta=$ & $1$ & $0.9$ & $0.8$ & $0.7$ & $0.65$ & $0.6$\\
\midrule
predicted sd of $\sigma_j$ ($\times10^{-4}$) & $25.1$ & $14.6$ & $9.1$ & $5.2$ & $3.2$ & $1.84$\\
particle filter, mean over six zeros & $27.1$ & $15.3$ & $9.4$ & $5.3$ & $3.3$ & $1.84$\\
\bottomrule
\end{tabular}
\caption{Posterior standard deviation of $\sigma_j$ for the six zeros of $\zeta$ in $[78,94]$, against the first order Fisher information formula ($\varepsilon=0.01$, $\Delta v=0.04$). The predicted row is computed (\texttt{code/cf\_info.py}); the filter row is the author's run.}
\label{cf:tab:info}
\end{table}

When the functional equation is imposed, a zero on the line is compared with a symmetric pair and the information is second order (Proposition~\ref{cf:prop:detect}): the relevant quantity is $\sum_la_l^{-5}$ and the resolvable $\delta$ scales like $a_{\min}^{5/4}$ (heuristic). Both formulas say the same thing: information comes from the levels nearest the line. The levels $\theta\ge0.9$ supply $1.6\%$ of $\sum a^{-3}$ and $0.1\%$ of $\sum a^{-5}$ in this schedule.

\begin{problem}\label{cf:op:pairres}
What is the best possible resolution of the pair parameter $\eta=\delta^2$ from data at levels $\theta\ge\theta_{\min}$, and does the particle filter with the pair parametrisation of Section~\ref{es:sec:singular} reach it? (open)
\end{problem}

\section{Optimising the filter with generative Bayesian computation}\label{cf:sec:gbc}

\subsection{Contour moments}\label{cf:sec:moments}

Before learning a summary, we recall what exact contour data give.

\begin{proposition}[contour moments]\label{cf:prop:moments}
Let $R=[\theta_0,\theta_1]\times[T_1,T_2]$ with $\xi\ne0$ on $\partial R$, let $c\in\mathbb C$ and
\[
M_k(R)=\frac1{2\pi i}\oint_{\partial R}(s-c)^k\,\frac{\xi'}{\xi}(s)\,ds\qquad(k\ge0).
\]
Then $M_k(R)=\sum_{\rho\in R}(\rho-c)^k$, with multiplicity. In particular $K=M_0(R)$ is the number of zeros in $R$, and $M_1(R),\dots,M_K(R)$ determine them: with $e_0=1$ and $je_j=\sum_{r=1}^j(-1)^{r-1}e_{j-r}M_r(R)$ for $1\le j\le K$, the numbers $\rho-c$, $\rho\in R$, are the roots of $z^K-e_1z^{K-1}+\dots+(-1)^Ke_K$.
\end{proposition}

\begin{proof}
The function $(s-c)^k\xi'(s)/\xi(s)$ is meromorphic inside $R$, with a simple pole at each zero $\rho$ of multiplicity $m$, of residue $m(\rho-c)^k$; the residue theorem gives the first claim. The power sums $M_1,\dots,M_K$ of the $K$ numbers $\rho-c$ determine their elementary symmetric functions $e_1,\dots,e_K$ by Newton's identities, which is the recursion displayed, and these are the coefficients of the monic polynomial with those roots.
\end{proof}

\begin{remark}
This is the method of Delves and Lyness \cite{DL67}. The moments use $\xi'/\xi$ itself on the whole of $\partial R$, horizontal edges included, at abscissae at least $\theta_0$; with exact contour integrals they give the zeros in $R$ exactly. They are not functions of the grid values used by the filter, so they are not a sufficient statistic of its data. The real part of an analytic function on the two vertical edges does not determine it: with $a=\theta_0$, $b=\theta_1$ and $T_1<t_0<T_2$, the entire function $g(s)=\cos\bigl(\pi(s-a-it_0)/(b-a)\bigr)$ is real and nonzero on both vertical lines, $g'/g$ is purely imaginary there, and $g$ vanishes at $\frac{a+b}2+it_0\in R$; so $\xi$ and $\xi g$ have the same density on both vertical edges and different zeros in $R$. Within the class of functions like $\xi$, of order one with zeros in the strip, the density on one segment of the line $\Re s=1$ does determine $\xi'/\xi$ (Proposition~\ref{cf:prop:determine}(a)), but only through analytic continuation, which no finite grid of values carries.
\end{remark}

\begin{table}[h]
\centering
\small\setlength{\tabcolsep}{4pt}\begin{tabular}{p{0.9in}cp{0.95in}p{1.4in}}
\toprule
function & rectangle & $M_0$ & recovered zeros right of $\theta_0$\\
\midrule
$\zeta$ & $[0.6,3]\times[80,92]$ & $0$ (radius $3\times10^{-48}$) & none\\
$\zeta$ & $[0.55,3]\times[60,120]$ & $0$ (radius $7\times10^{-46}$) & none\\
Davenport and Heilbronn & $[0.6,3]\times[80,92]$ & $1$ (radius $3\times10^{-48}$) & $0.808517182457+\allowbreak85.699348485378\,i$\\
\bottomrule
\end{tabular}
\caption{Contour moments, computed in ball arithmetic at $128$ and $192$ bits (\texttt{code/cf\_moments.py}, \texttt{data/cf\_moments.json}); the two precisions agree, and the integrals are enclosed, so that each $M_0$ is an enclosure containing one integer. The note printed $7\times10^{-22}$ and $0.006$ for the first two entries, floating point quadrature residuals (the second from zeros at distance $0.05$ from the edge); the exact values are $0$. The zero of the last row is the Newton root with $K=1$ and $M_1$ as an interval.}
\label{cf:tab:dl}
\end{table}

So in the noise free limit the deconvolution in a window is computed by a handful of contour integrals, to twelve digits for the zero of the last row. What makes the problem hard is not the window; it is that there are infinitely many windows. The role of GBC is to make each window nearly free, with its uncertainty quantified, when the data are a grid of noisy values rather than exact contour integrals.

\subsection{GBC: learned summary and quantile network}

Following Polson and Sokolov \cite{PS23}, and the semi automatic summaries of Fearnhead and Prangle \cite{FP12}, we simulate from the joint law of parameters and data, learn the conditional mean as a summary, and then learn the posterior quantile function given the summary (the book's account is Section~\ref{lo:sec:gbc}).

\paragraph{Simulator.} A window of half width $4$ around a centre $T$; ordinates from a rigid point process (Gamma$(4)$ spacings, density uniform on $[0.35,0.8]$ per unit, covering the Davenport and Heilbronn function and $\zeta$ at heights up to a few hundred; the density of $\zeta$ is $0.35$ at height $56.8$ and lower below, so the windows of $\zeta$ centred at $30$ to $55$ lie outside the training range); with probability $\frac12$ one zero in the central interval $\lvert\gamma-T\rvert\le3$ is replaced by a symmetric pair at distance $\delta$, mostly uniform on $[0.02,0.45]$; with probability $0.2$ a decoy pair outside the central interval; a random quadratic background proportional to $\theta-\frac12$; observations $u_\theta$ on $161$ points at $\theta\in\{1,0.85,0.7,0.6\}$ with error $0.003$. The target is $\delta^\ast=\max\{\lvert\sigma_\rho-\frac12\rvert:\lvert\gamma_\rho-T\rvert\le3\}$, equal to $0$ when the central zeros lie on the line.

\paragraph{Stage one: summary.} A network $S$ (two hidden layers, $256$ and $128$ units; the code's input has $652$ entries for the $644$ grid values, and the note does not say what the other eight are) is trained by squared error to $S(y)\approx\mathbb E[(\delta^\ast,\,\mathbf 1\{\delta^\ast>0\},\,\gamma^\ast)\mid y]$, with $\gamma^\ast$ the ordinate of the displaced zero. In the code the location loss is masked by the indicator that a displaced pair is present, so the third output estimates the location given that a pair is present, $\mathbb E[A\gamma^\ast\mid y]/\mathbb E[A\mid y]$ with $A=\mathbf 1\{\delta^\ast>0\}$, and the second network receives the three outputs together with eight hand built minima of the data, an eleven dimensional conditioning vector. By the argument of Fearnhead and Prangle the conditional mean is the optimal summary for estimating these targets under quadratic loss, when the network equals it.

\paragraph{Stage two: posterior by noise outsourcing.} A second network $Q(\tau,S(y))$ is trained by the pinball loss $\rho_\tau(\delta^\ast-Q)$ with $\tau\sim U(0,1)$ as an input. By the noise outsourcing theorem \cite{Kallenberg02}, $\delta^\ast\overset d=Q(U,S(y))$ with $U$ uniform and independent, for the exact conditional quantile function given $S(y)$, so posterior draws cost one forward pass each; the network approximates that function, and the posterior given $S(y)$ is the posterior given $y$ only if $S$ is sufficient (as in Section~\ref{lo:sec:gbc}). The network is not constrained to be monotone in $\tau$.

\paragraph{Training and calibration.} With $120{,}000$ training and $20{,}000$ test simulations, training took $111$ seconds on two cores. The root mean square error of the first output of $S$ for $\delta^\ast$ is $0.016$. Calling a window flagged when the posterior median exceeds $0.03$ gives a true positive rate of $0.978$ and a false positive rate of $0.002$; detection is essentially perfect for $\delta^\ast\ge0.05$ and falls to $0.67$ for $\delta^\ast\in[0.02,0.05]$ (computed by the author). Coverage of the posterior quantiles is in Table~\ref{cf:tab:cov}.

\begin{table}[h]
\centering
\begin{tabular}{lccccc}
\toprule
nominal $\tau$ & $0.05$ & $0.25$ & $0.50$ & $0.75$ & $0.95$\\
empirical & $0.049$ & $0.182$ & $0.467$ & $0.725$ & $0.917$\\
\bottomrule
\end{tabular}
\caption{Coverage of the GBC posterior quantiles of $\delta^\ast$ on held out simulations (computed by the author). The middle quantiles undercover, an approximation and calibration error of the quantile network; at an atom exact quantiles overcover. The upper tail is slightly short and is the place to add conformal calibration.}
\label{cf:tab:cov}
\end{table}

\subsection{Screening real windows}

We applied the trained networks to $35$ windows centred at $T=30,35,\dots,200$ for each function. Each window costs $644$ evaluations of $\xi'/\xi$ (of $\Lambda'/\Lambda$ for the Davenport and Heilbronn function) and then $6.7$ milliseconds of network time. Table~\ref{cf:tab:scan} lists the flagged windows.

\begin{table}[h]
\centering
\small\setlength{\tabcolsep}{4pt}\begin{tabular}{lcccc}
\toprule
window centre & GBC location & true zero & GBC median $\delta^\ast$ & true $\delta$\\
\midrule
$85$ & $85.68$ & $0.8085+85.6993\,i$ & $0.314$ & $0.3085$\\
$115$ & $114.08$ & $0.6508+114.1633\,i$ & $0.143$ & $0.1508$\\
$165$ & $166.47$ & $0.5744+166.4793\,i$ & $0.086$ & $0.0744$\\
$175$ & $176.72$ & $0.7243+176.7025\,i$ & $0.211$ & $0.2243$\\
$180$ & $176.58$ & same zero, at the edge & $0.361$ & $0.2243$\\
\bottomrule
\end{tabular}
\caption{Davenport and Heilbronn windows flagged by GBC (computed by the author; the zeros are those of Section~\ref{cf:sec:test}). The four zeros off the line below height $200$ are flagged, and located to within $0.1$ in the four windows that target them. The window at $180$ sees the zero at $176.70$ near its edge, outside its target interval $[177,183]$, and overstates $\delta$: a boundary false alarm, a boundary effect of the simulator. No other window of either function is flagged; for all $35$ windows of $\zeta$ the posterior median is $0$ and the $95\%$ quantile is at most $0.007$.}
\label{cf:tab:scan}
\end{table}

\subsection{GBC inside the particle filter}\label{cf:sec:gbcpf}

The GBC posterior is fast but approximate; the particle filter targets the model posterior but is slow. We combine them in the standard way, using the learned posterior as the starting distribution of the particle system.
\begin{enumerate}
\item Screen every window with GBC.
\item In a flagged window, build a proposal $q$: ordinates near the peaks of $u_{\theta_n}$ at the lowest level for zeros on the line, a pair near each negative dip with $\delta$ drawn from the GBC posterior by noise outsourcing, and diffuse components for the rest.
\item Run the tempered sampler from $q$ to the posterior, through the bridge $q^{1-\lambda}(\pi_0L)^\lambda$, whose potentials are powers of $\pi_0L/q$ and so involve $\log\pi_0-\log q$, with the same resample and move steps as in Section~\ref{cf:sec:pf}.
\end{enumerate}
In practice step 3 works far better when the proposal is first moved to the posterior mode by a Gauss Newton fit started from the GBC output, with $q$ the Gaussian with four times the inverse Gauss Newton Hessian as covariance. The results:

\begin{table}[h]
\centering
\small\begin{tabular}{p{5.2cm}cccc}
\toprule
sampler & window & steps & evaluations & seconds\\
\midrule
plain filter, uniform prior, $\zeta$ & $[78,94]$ & many & $4.15\times10^6$ & minutes\\
GBC, Gauss Newton, filter, $\zeta$ & $[80,92]$ & $7$ & $3.6\times10^4$ & $3.6$\\
plain filter, uniform prior, Davenport and Heilbronn & $[78,94]$ & stalled & n/a & stopped\\
GBC, Gauss Newton, filter, Davenport and Heilbronn & $[80,92]$ & $1020$ & $5.1\times10^6$ & $645$\\
\bottomrule
\end{tabular}
\caption{Costs of reported runs (computed by the author). The two runs for $\zeta$ differ in window ($[78,94]$ with six zeros against $[80,92]$ with four), in zeros and in data, so their ratio of about $115$ in evaluations is a comparison of two different runs and not a matched speedup. For Davenport and Heilbronn the GBC initialised filter succeeds where the plain filter stalls (at tempering exponent $\lambda\approx0.003$), but the posterior is strongly non Gaussian because the level $\theta=0.8$ lies within $0.0085$ of $\sigma=0.8085$, so the bridge from the Gaussian proposal still needs many steps.}
\label{cf:tab:cost}
\end{table}

For $\zeta$ the four zeros inside $[80,92]$ are recovered as before, with $95\%$ intervals for $\sigma_j$ of half width $4\times10^{-4}$ around $\frac12$. For the Davenport and Heilbronn function the filter finds the pair at $\gamma=85.6993$ with $\sigma=0.80851$ (interval $[0.80851,0.80852]$) and $\sigma=0.18994$ (interval $[0.18867,0.19084]$). The first agrees with the true $0.808517$ to $10^{-5}$. The second misses the true $0.191483$ by $0.0015$, and the intervals for several zeros on the line exclude $\frac12$ by up to $6\times10^{-4}$. Both are symptoms of misspecification: with $12$ zeros in a window of length $18$ the cubic background is too crude for a tolerance of $\varepsilon=0.01$, and the posterior is overconfident. A larger zero window, or a background built from the Riemann von Mangoldt density, should help; the last paragraph of the next subsection confirms the first for these data.

\subsection{Three computational questions}\label{cf:sec:questions}

\paragraph{Sine kernel statistics in the simulator.} We replaced the rigid Gamma$(4)$ spacing of the GBC simulator by eigenphases of Haar unitary matrices (the circular unitary ensemble), rescaled to the local density, whose local statistics approximate the sine kernel correlations conjectured for the zeros of $\zeta$ (the sine process is the limit as the size grows, and the matrix size is not stated in the note, heuristic); the empirical spacing variance is $0.16$ against $0.25$ before. Retraining both networks took $285$ seconds. On held out simulations of the new simulator the $95\%$ quantile has coverage $0.952$, against $0.917$ for the rigid networks on theirs. On the real windows the screen again flags exactly the four Davenport and Heilbronn zeros, plus the edge window at $180$, and nothing for $\zeta$ (largest posterior median $0.001$). The estimates of $\delta$ improve about threefold (Table~\ref{cf:tab:gue}, computed by the author).

\begin{table}[h]
\centering
\small\begin{tabular}{lcccc}
\toprule
zero & \shortstack{$0.8085$\\$+85.70\,i$} & \shortstack{$0.6508$\\$+114.16\,i$} & \shortstack{$0.5744$\\$+166.48\,i$} & \shortstack{$0.7243$\\$+176.70\,i$}\\
\midrule
true $\delta$ & $0.3085$ & $0.1508$ & $0.0744$ & $0.2243$\\
rigid simulator & $0.314$ & $0.143$ & $0.086$ & $0.211$\\
sine kernel simulator & $0.312$ & $0.149$ & $0.079$ & $0.228$\\
\bottomrule
\end{tabular}
\caption{GBC posterior medians of $\delta$. Mean absolute error $0.0096$ with the rigid simulator and $0.0034$ with the sine kernel simulator.}
\label{cf:tab:gue}
\end{table}

\paragraph{Conformal calibration.} For the rigid networks, split conformal calibration on $10{,}000$ fresh simulations shifts the $95\%$ upper bound by an order statistic of the calibration residuals, $0.0015$, and raises its coverage on $20{,}000$ test simulations from $0.915$ to $0.948$. Conditional on an off line pair being present, coverage rises only from $0.885$ to $0.906$: marginal conformal calibration does not repair conditional coverage. A Mondrian version partitions by a label that must be observed at prediction time; the indicator $\mathbf 1\{\delta^\ast>0\}$ is not, so the partition would have to use the network's own flag (open). For the sine kernel networks the shift is $-0.0002$ and no correction is needed.

\paragraph{Removing the misspecification.} Section~\ref{cf:sec:gbcpf} reported that the filter missed the partner zero, with $\sigma=0.19148$ at height $85.6993$, by $0.0015$ and was overconfident on the zeros on the line. We attributed this to too narrow a zero window. Widening it to $[74,98]$ with $K=17$ zeros, counted by the argument principle, and data on $[78,94]$ at the five levels other than $\theta=0.8$, the Gauss Newton fit started from the GBC output converged in $82$ evaluations and $4.7$ seconds, with essentially zero residual. With Laplace standard errors it gives
\[
\sigma=0.80852\pm0.00015,\qquad 1-\sigma=0.19149\pm0.00129,\qquad \gamma=85.6993\pm0.0001,
\]
against the true $0.808517$, $0.191483$ and $85.699348$. Every zero on the line in the data window comes out as $0.50000\pm0.00019$. The wider window fit removes the discrepancy of the narrower fit for these data. The six components outside the data window are unidentified and drift to the boundary of the prior box, as nuisance parameters should. They are also why the full tempered sampler on this wider window did not finish in forty minutes: the posterior in those coordinates is flat and truncated, far from the Gaussian proposal. Integrating the unobserved zeros into the background, through the Riemann von Mangoldt density, rather than carrying them as particles, is a remedy that has not been tried.
